% Generated draft: requires mathematical and source review.
\documentclass[11pt]{article}
\usepackage[margin=25mm]{geometry}
\usepackage{fontspec}
\setmainfont{Latin Modern Roman}
\usepackage{amsmath,amssymb,mathtools,unicode-math}
\setmathfont{Latin Modern Math}
\usepackage{enumitem,needspace,xurl,hyperref}
\setlength{\parindent}{0pt}
\setlength{\parskip}{5pt}
\setlength{\emergencystretch}{3em}
\newcommand{\N}{\mathbb{N}}
\newcommand{\Z}{\mathbb{Z}}
\newcommand{\Q}{\mathbb{Q}}
\newcommand{\R}{\mathbb{R}}
\newcommand{\C}{\mathbb{C}}
\newcommand{\F}{\mathbb{F}}
\newcommand{\A}{\mathbb{A}}
\newcommand{\PP}{\mathbb{P}}
\newcommand{\E}{\mathbb{E}}
\newcommand{\eps}{\varepsilon}
\newcommand{\dd}{\,\mathrm{d}}
\begin{document}
\section*{UnsolvedMath ID: 700987}
\textbf{Source problem code:} catalogue entry
\subsection{Definitions and mathematical statement}
Let \(G\) be a finite group. Its integral group ring is
\[
\Z[G]=\left\{\sum_{g\in G}a_g g:a_g\in\Z\right\},
\]
with multiplication extending the multiplication of \(G\) distributively. Its centre is
\[
Z(\Z[G])=\{z\in\Z[G]:zx=xz\text{ for every }x\in\Z[G]\}.
\]
Let \(\mathcal C(G)\) denote the set of conjugacy classes of \(G\). For \(C\in\mathcal C(G)\), define its class sum by
\[
K_C=\sum_{g\in C}g.
\]
The elements \(K_C\), for \(C\in\mathcal C(G)\), form a \(\Z\)-basis of \(Z(\Z[G])\).

Interpret an automorphism as a unital ring automorphism of \(Z(\Z[G])\); such an automorphism fixes the embedded integers and is \(\Z\)-linear. Interpret a monomial permutation as an action whose matrix in the class-sum basis has exactly one nonzero entry in every row and every column. Since this matrix is invertible over \(\Z\), those entries must lie in \(\{1,-1\}\).

\textbf{Statement requested by the source, under these conventions.} For every finite group \(G\) and every ring automorphism
\[
\alpha:Z(\Z[G])\longrightarrow Z(\Z[G]),
\]
there exist a bijection \(\pi:\mathcal C(G)\longrightarrow\mathcal C(G)\) and a function \(\varepsilon:\mathcal C(G)\longrightarrow\{1,-1\}\) such that
\[
\alpha(K_C)=\varepsilon(C)K_{\pi(C)}
\qquad\text{for every }C\in\mathcal C(G).
\]
The supplied record does not define ``monomial permutation'' or explicitly specify the category of automorphisms. The displayed formulation uses the standard ring and integral monomial-matrix interpretations; these conventions require confirmation against the cited source.
\subsection{Short English statement}
For every finite group, the requested assertion is that each ring automorphism of the centre of its integral group ring permutes the conjugacy-class sums, allowing a sign on each sum. This uses the signed-permutation interpretation of ``monomial permutation,'' which the supplied record leaves undefined.
\subsection{Sources}
\begin{itemize}
\item Local UnsolvedMath registry: permanent ID \(700987\), exact source code \texttt{catalogue entry}, title \emph{Kourovka Notebook Problem 14.2}; supplied upstream revision \texttt{d371ee917ac9647fbb06cd6fc529deec6f3dda82}.
\item Supplied source attribution: \emph{The Kourovka Notebook (Unsolved Problems in Group Theory)}, problem 14.2, 14th issue (1999); the record also cites the 21st issue edition (2026), arXiv:1401.0300v48. The collection description names E. I. Khukhro and V. D. Mazurov as editors. These bibliographic details are supplied and have not been independently verified.
\item The record names R. Zh. Aleev as proposer, with proposed year 1999. The statement itself carries the attribution ``(S. D. Berman).'' Both attributions are retained; their precise relationship is not explained in the supplied record.
\item The local registry labels the problem \texttt{open}. This is a supplied status label, not an independently verified assessment of the literature.
\end{itemize}
\end{document}
